\DocumentMetadata{
  pdfversion=2.0,
  pdfstandard=ua-2,
  lang=en-US,
  tagging=on,
  tagging-setup={math/setup=mathml-SE,tabsorder=structure}
}
\documentclass[11pt,letterpaper]{article}
\usepackage[margin=0.68in,includefoot]{geometry}
\usepackage{newcomputermodern}
\usepackage{amsmath,amscd,mathtools}
\usepackage{tikz,adjustbox}
\providecommand{\square}{\mdlgwhtsquare}
\usepackage[hidelinks]{hyperref}
\hypersetup{
  pdftitle={Analysis First-Year Exam, August 2015, Part 2},
  pdfauthor={University of Florida Department of Mathematics},
  pdfsubject={Graduate mathematics examination},
  pdfdisplaydoctitle=true,
  bookmarksopen=true
}
\setcounter{secnumdepth}{0}
\setlength{\parindent}{0pt}
\setlength{\parskip}{0.28em}
\setlength{\emergencystretch}{3em}
\setlength{\leftmargini}{2.15em}
\setlength{\leftmarginii}{2.65em}
\setlength{\leftmarginiii}{3em}
\renewcommand{\labelenumi}{\textbf{\theenumi.}}
\pagestyle{empty}
\raggedbottom
\allowdisplaybreaks
\newenvironment{examproblems}[1]{%
  \begin{enumerate}%
  \setcounter{enumi}{#1}%
  \setlength{\topsep}{0.35em}%
  \setlength{\partopsep}{0pt}%
  \setlength{\itemsep}{0.48em}%
  \setlength{\parsep}{0.18em}%
}{\end{enumerate}}
\newenvironment{examsubparts}{%
  \begin{enumerate}%
  \setlength{\topsep}{0.18em}%
  \setlength{\partopsep}{0pt}%
  \setlength{\itemsep}{0.16em}%
  \setlength{\parsep}{0.1em}%
}{\end{enumerate}}
\newenvironment{examinstructions}{%
  \par\begingroup\itshape
}{%
  \par\endgroup\vspace{0.18em}
}
\makeatletter
\renewcommand\section{\@startsection{section}{1}{0pt}%
  {-1.25ex plus -.25ex minus -.1ex}%
  {0.55ex plus .1ex}%
  {\normalfont\large\bfseries}}
\renewcommand{\@maketitle}{%
  \newpage\null\vskip -1em%
  {\footnotesize\bfseries
    \href{https://www.ufl.edu/}{University of Florida}, \href{https://math.ufl.edu/}{Department of Mathematics}\par}%
  \vskip -0.5em%
  \tagstructbegin{tag=Title}%
    {\LARGE\bfseries\href{https://gma.math.ufl.edu/exams/analysis-first-year/}{Analysis First-Year Exam}, August 2015, Part 2\par}%
  \tagstructend%
  \vskip 0.25em%
  \tagmcbegin{artifact}%
    \hrule height 1pt%
  \tagmcend%
  \vskip 1em%
}
\makeatother
\title{Analysis First-Year Exam, August 2015, Part 2}
\author{}
\date{}
\begin{document}
\maketitle
\thispagestyle{empty}
\begin{examinstructions}
Do exactly 2 problems from Part A and 2 problems from Part B. Work must be presented in a neat and logical fashion in order to receive credit. Do not leave any gaps. When a theorem is used in a proof it must be precisely stated.
\end{examinstructions}
\section{Part A}
\begin{examproblems}{0}
\item
For \(n \ge 1\) and \(0<t<1\), let \(f_n(t)=t^n\). Prove that:
\begin{examsubparts}
\item[\textnormal{a)}]
\((f_n)\) converges uniformly on each compact subset of \((0,1)\),
\item[\textnormal{b)}]
\((f_n)\) does not converge uniformly on \((0,1)\).
\end{examsubparts}
\item
Suppose \(\mathcal F\) is an equicontinuous family of real-valued functions on \([0,1]\), such that 
\[
\sup_{f\in\mathcal F}|f(0)|=M<+\infty.
\]
 Prove that \(\mathcal F\) is uniformly bounded (that is, that \(\sup_{f\in\mathcal F}\sup_{x\in[0,1]}|f(x)|<+\infty\)).
\item
Prove that every monotone function \(f:[a,b]\to\mathbb R\) is Riemann integrable.
\end{examproblems}
\section{Part B}
\begin{examproblems}{0}
\item
State Fatou’s lemma and use it to prove the Dominated Convergence Theorem.
\item
Let \((f_n)\) be a sequence of real-valued measurable functions on a measurable space \((X,\mathcal M)\). Prove that each of the following subsets of~\(X\) is measurable:
\begin{examsubparts}
\item[\textnormal{a)}]
\(\{x\in X:\text{ the sequence }(f_n(x))\text{ is unbounded}\}\).
\item[\textnormal{b)}]
\(\{x\in X:\text{ the sequence }(f_n(x))\text{ is strictly increasing}\}\).
\end{examsubparts}
\item
Let \(g:[0,1]\times[0,1]\to\mathbb R\) be a function with the following properties. Prove that the function~\(h\) defined by 
\[
h(x)=\int_0^1 g(x,t)\,dt
\]
 is continuous on \([0,1]\).
\begin{examsubparts}
\item[\textnormal{i)}]
\(|g(x,t)|\le 1\) for all~\(x\) and~\(t\),
\item[\textnormal{ii)}]
for each~\(t\), the function \(x\mapsto g(x,t)\) is continuous on \([0,1]\), and
\item[\textnormal{iii)}]
for each~\(x\), the function \(t\mapsto g(x,t)\) is continuous on \([0,1]\).
\end{examsubparts}
\end{examproblems}
\end{document}
